% =====================================================================
\subsection{Limitations and Future Work}
\label{sec:limits}
% =====================================================================

\paragraph{Current limitations.}
Level~0 and Level~1 canonicalization do not capture full definitional
equality; the gap between Level~1 and the ideal is large and may be
fundamentally unclosable.  Level~1 addresses depend on Lean's
\texttt{whnf} behavior, which can change across toolchain versions;
Level~0 addresses are stable within a major Lean version, while Level~1
should be considered best-effort within a given toolchain.  Lean's
native \texttt{whnf} (implemented in C) lacks heartbeat checks and
crashes on large environments; a micro-batching workaround is effective
but inelegant (Appendix~\ref{app:engineering}).  Applying \texttt{whnf}
to definition values causes process-level memory growth that crashes
regardless of batch boundaries, limiting content-addressing of
definitions (as opposed to theorems, where proof irrelevance makes value
hashing unnecessary).  Finally, Lean~4's elaboration pipeline does not
support user-defined fallbacks for unresolved names; our resolution
mechanism requires explicit syntax (\texttt{resolve\%}), preventing
seamless integration.

\paragraph{Future work.}
\phantomsection\label{sec:future}
Several directions are worth pursuing.  \emph{Resolution elaborators}:
implement \texttt{use}, \texttt{use!}, and \texttt{resolve\%} as Lean~4
elaborators that query the registry, resolve by content address, and
dynamically import modules.  The intended syntax (not yet implemented):

\begin{lstlisting}[caption={Prospective resolution syntax (not yet implemented)}]
-- resolve% looks up a content address in the registry and
-- returns the name of a matching declaration in scope
example : forall a b : Nat, a + b = b + a :=
  resolve% "a1b2c3d4e5f6..."   -- resolves to Nat.add_comm
\end{lstlisting}

\emph{Ecosystem-wide indexing}: apply
content-addressing to the entirety of Mathlib to produce a global index
of equivalence classes, enabling large-scale duplicate detection,
library health metrics, and structured navigation.
\emph{Open-problem workflows}: develop the open-point and
conditional-result mechanism (Section~\ref{sec:results}) into a robust
workflow, addressing false conjectures, \texttt{sorry}-based vs.\
axiomatic assumptions, and cross-project status curation.
\emph{AI-assisted proof search}: content addresses provide a stable,
type-based index for machine learning models to retrieve relevant lemmas
by querying the content address of a goal type.
\emph{Cross-system alignment}: explore content-addressing in other proof
assistants (Coq, Isabelle, Agda) and investigate whether addresses can
be made comparable via shared canonicalization.
\emph{Intermediate canonicalization levels}: investigate canonicalization
at semireducible transparency or with selective instance unfolding.
\emph{Formal verification}: prove (in Lean itself) that the
canonicalization and serialization stages preserve the intended
equivalence relation.

\moqian{CA + ZKP reduces verification from O(compile entire dependency chain) to O(1). Critical for AI-generated proofs at scale.}

\moqian{CA shortens proof writing: hash as stable ID, readable name as alias. DNS model---IP is identity, domain name is human-friendly.}

\paragraph{Engineering Obstacles}
\label{app:engineering}

A practical obstacle to Level~1 canonicalization at scale: Lean's native
\texttt{whnf} (implemented in C) contains no heartbeat or memory
checks.  When processing large environments (${\sim}58{,}000$
declarations in Mathlib modules), two failure modes arise.  First,
\emph{MetaM state accumulation}: a single \texttt{MetaM} session
crashes after ${\sim}25{,}000$ \texttt{whnf} calls due to cumulative
growth of caches, discriminant trees, and instance tables; the crash
point depends on call count, not on which specific declarations are
processed.  Second, \emph{process-level memory growth}: applying
\texttt{whnf} to definition values causes monotonic RSS growth that
eventually crashes regardless of batch boundaries.

\paragraph{Micro-batching.}  Our workaround is to process declarations
in micro-batches of~10, each in a fresh \texttt{MetaM} session.  Types
receive Level~1 normalization; values are always canonicalized with
Level~0 only (no \texttt{whnf}).  For theorems, the value hash is the
sentinel \texttt{PROOF\_IRRELEVANT} (proof irrelevance).  The result is
an asymmetric strategy: \textbf{types get L1, values get L0}.  Hash
computation is split into two phases (canonicalize+hash vs.\
pretty-print) in separate \texttt{MetaM} sessions, with a flush every
5{,}000 declarations to bound live \texttt{Expr} objects in memory.

An upstream fix---heartbeat checks and resource guards in Lean's native
\texttt{whnf} implementation---would eliminate the need for
micro-batching entirely.

% =====================================================================
